% GENERATED FILE. Edit canonical lessons, then run npm run generate:books.
% canonical-source-hash: fnv1a64:a974a2e774383857
% canonical-lessons: JA-C371-joudan, JA-C371-sankou, JA-C371-waruguchi, JA-C371-tereru, JA-C371-shaberu, JA-R371-review-367-371, JA-R371-review-final-371

\chapter{Seriously or Joking?}
\label{ch:ja-seriously-or-joking}
\hlchaptermodality{371}

\begin{chapteropening}
\textbf{By the end of this chapter:} \emph{I can tell a joke from a serious remark and say what would follow.}

Complete the last lesson of chapter 371: I can tell a joke from a serious remark and say what would follow.
\end{chapteropening}

\section[\texorpdfstring{\ja{じょうだん} (jōdan)}{じょうだん (jōdan)}]{\ja{じょうだん} (jōdan) — a joke}

\label{lesson:JA-C371-joudan}



\begin{quote}
Before the new one: say the Japanese for roughly, then the Japanese for perfectly.
\end{quote}

\subsection*{You'll want to know: \ja{じょうだん}}
\textbf{\ja{じょうだん}} — \emph{jōdan} — \textquotedblleft{}a joke\textquotedblright{}.

\begin{quote}
\textbf{\ja{じょうだんを} \ja{言}\ja{ったら、みんな} \ja{わらうでしょう。}} — \emph{jōdan o ittara, minna warau deshō} — \textquotedblleft{}If you tell a joke, everyone will probably laugh.\textquotedblright{}
\end{quote}

\begin{cousinweb}
\textbf{\ja{じょうだん}} (\emph{jōdan}) is Sino-Japanese (kango): two Chinese characters read with their on'yomi, readings borrowed from Chinese. They are \emph{jō}, \textquotedblleft{}superfluous\textquotedblright{} and \emph{dan}, \textquotedblleft{}discuss\textquotedblright{} (Middle Chinese \emph{dam}) (KANJIDIC2; Middle Chinese in Baxter's transcription).
\end{cousinweb}

\begin{grammarlens}[title={What you've built}]
One more word for saying what happens if.
\end{grammarlens}

\subsection*{Your turn}
\begin{itemize}
\raggedright
  \item \emph{Say it:} \emph{jōdan}
  \item \emph{Say it:} \emph{jōdan}, once more
  \item \emph{Say it:} \emph{jōdan o ittara, minna warau deshō} — If you tell a joke, everyone will probably laugh.
  \item \emph{Recall:} say the Japanese for roughly, then the Japanese for perfectly, then say \emph{jōdan} again
\end{itemize}

\begin{tcolorbox}[breakable,colback=teal!4,colframe=teal!35!black,title={Before you move on}]
What does \ja{じょうだん} mean? (\textquotedblleft{}A joke\textquotedblright{}.) Say it once more.
\end{tcolorbox}
\section[\texorpdfstring{\ja{さんこう} (sankō)}{さんこう (sankō)}]{\ja{さんこう} (sankō) — a reference; something to learn from}

\label{lesson:JA-C371-sankou}



\begin{quote}
Before the new one: say the Japanese for perfectly, then the Japanese for a joke.
\end{quote}

\subsection*{You'll want to know: \ja{さんこう}}
\textbf{\ja{さんこう}} — \emph{sankō} — \textquotedblleft{}a reference; something to learn from\textquotedblright{}.

\begin{quote}
\textbf{\ja{この} \ja{ほんを} \ja{よんだら、さんこうに} \ja{なるでしょう。}} — \emph{kono hon o yondara, sankō ni naru deshō} — \textquotedblleft{}If you read this book, you will probably find it useful.\textquotedblright{}
\end{quote}

\begin{cousinweb}
\textbf{\ja{さんこう}} (\emph{sankō}) is Sino-Japanese (kango): two Chinese characters read with their on'yomi, readings borrowed from Chinese. They are \emph{san}, \textquotedblleft{}participate\textquotedblright{} and \emph{kō}, \textquotedblleft{}consider\textquotedblright{} (Middle Chinese \emph{khawX}) (KANJIDIC2; Middle Chinese in Baxter's transcription).
\end{cousinweb}

\begin{grammarlens}[title={What you've built}]
One more word for saying what happens if.
\end{grammarlens}

\subsection*{Your turn}
\begin{itemize}
\raggedright
  \item \emph{Say it:} \emph{sankō}
  \item \emph{Say it:} \emph{sankō}, once more
  \item \emph{Say it:} \emph{kono hon o yondara, sankō ni naru deshō} — If you read this book, you will probably find it useful.
  \item \emph{Recall:} say the Japanese for perfectly, then the Japanese for a joke, then say \emph{sankō} again
\end{itemize}

\begin{tcolorbox}[breakable,colback=teal!4,colframe=teal!35!black,title={Before you move on}]
What does \ja{さんこう} mean? (\textquotedblleft{}A reference; something to learn from\textquotedblright{}.) Say it once more.
\end{tcolorbox}
\section[\texorpdfstring{\ja{わるぐち} (waruguchi)}{わるぐち (waruguchi)}]{\ja{わるぐち} (waruguchi) — speaking ill of someone; abuse}

\label{lesson:JA-C371-waruguchi}



\begin{quote}
Before the new one: say the Japanese for a joke, then the Japanese for a reference.
\end{quote}

\subsection*{You'll want to know: \ja{わるぐち}}
\textbf{\ja{わるぐち}} — \emph{waruguchi} — \textquotedblleft{}speaking ill of someone; abuse\textquotedblright{}.

\begin{quote}
\textbf{\ja{わるぐちを} \ja{言}\ja{ったら、ともだちが} \ja{おこるでしょう。}} — \emph{waruguchi o ittara, tomodachi ga okoru deshō} — \textquotedblleft{}If you speak ill of people, your friends will be angry.\textquotedblright{}
\end{quote}

\begin{cousinweb}
\textbf{\ja{わるぐち}} (\emph{waruguchi}) is a native Japanese word (wago), inherited rather than borrowed from Chinese: it is built from \emph{waru}, \textquotedblleft{}bad\textquotedblright{}, the stem of \emph{warui}, and \emph{guchi}, \textquotedblleft{}mouth\textquotedblright{}, the native \emph{kuchi} voiced, and its kanji were attached later as kun'yomi (JMdict, KANJIDIC2).
\end{cousinweb}

\begin{grammarlens}[title={What you've built}]
One more word for saying what happens if.
\end{grammarlens}

\subsection*{Your turn}
\begin{itemize}
\raggedright
  \item \emph{Say it:} \emph{waruguchi}
  \item \emph{Say it:} \emph{waruguchi}, once more
  \item \emph{Say it:} \emph{waruguchi o ittara, tomodachi ga okoru deshō} — If you speak ill of people, your friends will be angry.
  \item \emph{Recall:} say the Japanese for a joke, then the Japanese for a reference, then say \emph{waruguchi} again
\end{itemize}

\begin{tcolorbox}[breakable,colback=teal!4,colframe=teal!35!black,title={Before you move on}]
What does \ja{わるぐち} mean? (\textquotedblleft{}Speaking ill of someone; abuse\textquotedblright{}.) Say it once more.
\end{tcolorbox}
\section[\texorpdfstring{\ja{てれる} (tereru)}{てれる (tereru)}]{\ja{てれる} (tereru) — to be shy; to feel embarrassed}

\label{lesson:JA-C371-tereru}



\begin{quote}
Before the new one: say the Japanese for a reference, then the Japanese for speaking ill of someone.
\end{quote}

\subsection*{You'll want to know: \ja{てれる}}
\textbf{\ja{てれる}} — \emph{tereru} — \textquotedblleft{}to be shy; to feel embarrassed\textquotedblright{}.

\begin{quote}
\textbf{\ja{ほめたら、あには} \ja{てれるでしょう。}} — \emph{hometara, ani wa tereru deshō} — \textquotedblleft{}If you praise him, my brother will probably get embarrassed.\textquotedblright{}
\end{quote}

Polite \textbf{\ja{てれます}}, \emph{teremasu}; an ichidan verb. Past (polite): \textbf{\ja{あには} \ja{てれました。}}, \emph{ani wa teremashita}.

\begin{cousinweb}
\textbf{\ja{てれる}} (\emph{tereru}) is a native Japanese word (wago), inherited rather than borrowed from Chinese, and its older origin is not settled. Its kanji was attached to it later as a kun'yomi, a native reading given to a Chinese character (JMdict, KANJIDIC2).
\end{cousinweb}

\begin{grammarlens}[title={What you've built}]
One more word for saying what happens if.
\end{grammarlens}

\subsection*{Your turn}
\begin{itemize}
\raggedright
  \item \emph{Say it:} \emph{tereru}
  \item \emph{Say it:} \emph{tereru}, once more
  \item \emph{Say it:} \emph{hometara, ani wa tereru deshō} — If you praise him, my brother will probably get embarrassed.
  \item \emph{Recall:} say the Japanese for a reference, then the Japanese for speaking ill of someone, then say \emph{tereru} again
\end{itemize}

\begin{tcolorbox}[breakable,colback=teal!4,colframe=teal!35!black,title={Before you move on}]
What does \ja{てれる} mean? (\textquotedblleft{}To be shy; to feel embarrassed\textquotedblright{}.) Say it once more.
\end{tcolorbox}
\section[\texorpdfstring{\ja{しゃべる} (shaberu)}{しゃべる (shaberu)}]{\ja{しゃべる} (shaberu) — to chat; to talk}

\label{lesson:JA-C371-shaberu}



\begin{quote}
Before the new one: say the Japanese for speaking ill of someone, then the Japanese for to feel embarrassed.
\end{quote}

\subsection*{You'll want to know: \ja{しゃべる}}
\textbf{\ja{しゃべる}} — \emph{shaberu} — \textquotedblleft{}to chat; to talk\textquotedblright{}.

\begin{quote}
\textbf{\ja{じゅぎょうで} \ja{しゃべったら、めいわくです。}} — \emph{jugyō de shabettara, meiwaku desu} — \textquotedblleft{}Chatting in class is a nuisance to others.\textquotedblright{}
\end{quote}

Polite \textbf{\ja{しゃべります}}, \emph{shaberimasu}; a godan verb. Past (polite): \textbf{\ja{いもうとが} \ja{しゃべりました。}}, \emph{imōto ga shaberimashita}.

\begin{cousinweb}
\textbf{\ja{しゃべる}} (\emph{shaberu}) is a native Japanese word (wago), inherited rather than borrowed from Chinese, and its older origin is not settled. JMdict notes that it is usually written in kana (JMdict).
\end{cousinweb}

\begin{grammarlens}[title={What you've built}]
One more word for saying what happens if.
\end{grammarlens}

\subsection*{Your turn}
\begin{itemize}
\raggedright
  \item \emph{Say it:} \emph{shaberu}
  \item \emph{Say it:} \emph{shaberu}, once more
  \item \emph{Say it:} \emph{jugyō de shabettara, meiwaku desu} — Chatting in class is a nuisance to others.
  \item \emph{Recall:} say the Japanese for speaking ill of someone, then the Japanese for to feel embarrassed, then say \emph{shaberu} again
\end{itemize}

\begin{tcolorbox}[breakable,colback=teal!4,colframe=teal!35!black,title={Before you move on}]
What does \ja{しゃべる} mean? (\textquotedblleft{}To chat; to talk\textquotedblright{}.) Say it once more.
\end{tcolorbox}
\section[(dialogue)]{Looking back: from Being Clear to Seriously or Joking?}

\label{lesson:JA-R371-review-367-371}



\begin{quote}
Say the Japanese for clearly and in other words.
\end{quote}

\subsection*{Your turn}
Say these aloud:

\begin{itemize}
\raggedright
  \item clearly, in other words, enough, a degree, a proportion
  \item the whole, a section, both, one of a pair, to be one-sided
  \item essential, important, special, a demand, a policy
  \item to be settled, each other, the same, roughly, perfectly
  \item a joke, a reference, speaking ill of someone, to feel embarrassed, to chat
\end{itemize}

\begin{tcolorbox}[breakable,colback=teal!4,colframe=teal!35!black,title={Before you move on}]
Clearly? (\textbf{\ja{はっきり}}, \emph{hakkiri}.) In other words? (\textbf{\ja{つまり}}, \emph{tsumari}.) Enough? (\textbf{\ja{じゅうぶん}}, \emph{jūbun}.) A degree? (\textbf{\ja{ていど}}, \emph{teido}.) A proportion? (\textbf{\ja{わりあい}}, \emph{wariai}.) The whole? (\textbf{\ja{ぜんたい}}, \emph{zentai}.) A section? (\textbf{\ja{ぶぶん}}, \emph{bubun}.) Both? (\textbf{\ja{りょうほう}}, \emph{ryōhō}.) One of a pair? (\textbf{\ja{かたほう}}, \emph{katahō}.) To be one-sided? (\textbf{\ja{かたよる}}, \emph{katayoru}.) Essential? (\textbf{\ja{じゅうよう}}, \emph{jūyō}.) Important? (\textbf{\ja{だいじ}}, \emph{daiji}.) Special? (\textbf{\ja{とくべつ}}, \emph{tokubetsu}.) A demand? (\textbf{\ja{ようきゅう}}, \emph{yōkyū}.) A policy? (\textbf{\ja{ほうしん}}, \emph{hōshin}.) To be settled? (\textbf{\ja{まとまる}}, \emph{matomaru}.) Each other? (\textbf{\ja{たがいに}}, \emph{tagaini}.) The same? (\textbf{\ja{おなじ}}, \emph{onaji}.) Roughly? (\textbf{\ja{だいたい}}, \emph{daitai}.) Perfectly? (\textbf{\ja{ぴったり}}, \emph{pittari}.) A joke? (\textbf{\ja{じょうだん}}, \emph{jōdan}.) A reference? (\textbf{\ja{さんこう}}, \emph{sankō}.) Speaking ill of someone? (\textbf{\ja{わるぐち}}, \emph{waruguchi}.) To feel embarrassed? (\textbf{\ja{てれる}}, \emph{tereru}.) To chat? (\textbf{\ja{しゃべる}}, \emph{shaberu}.)
\end{tcolorbox}
\section[(dialogue)]{Seriously or Joking? — once more}

\label{lesson:JA-R371-review-final-371}



\begin{quote}
Say the Japanese for a joke and a reference.
\end{quote}

\subsection*{Your turn}
\begin{itemize}
\raggedright
  \item \emph{Say it:} a joke, a reference, speaking ill of someone, to feel embarrassed, to chat
\end{itemize}

\begin{tcolorbox}[breakable,colback=teal!4,colframe=teal!35!black,title={Before you move on}]
A joke? (\textbf{\ja{じょうだん}}, \emph{jōdan}.) A reference? (\textbf{\ja{さんこう}}, \emph{sankō}.) Speaking ill of someone? (\textbf{\ja{わるぐち}}, \emph{waruguchi}.) To feel embarrassed? (\textbf{\ja{てれる}}, \emph{tereru}.) To chat? (\textbf{\ja{しゃべる}}, \emph{shaberu}.)
\end{tcolorbox}
